\section{The common ordinary coefficient system and its arithmetic action}
\label{sec:common-ordinary-arithmetic-action}

We supply the finite-level arithmetic input to the common
patched module of Section~\ref{sec:sr-common-patched-module}.
The coefficient system and the base support ideal remain fixed
as auxiliary primes and auxiliary compact subgroups vary.
In particular, we do not replace the coefficient ring separately
by the image algebra on each auxiliary module.

Fix either of the actual rank-two or rank-three packets in
Sections~\ref{sec:seed-all-generic-primes} and
\ref{sec:arithmetic-frame-uniform-depth}.
Write \(F\) for the totally real field, \(L/F\) for its CM
extension, \(d=[F:\mathbb Q]\), and \(n\in\{2,3\}\).
Let \(S\) include the original level, the places above \(\ell\),
and the distinguished neat place \(v_0\).
The constructed fields satisfy
\[
 \ell\ne2,\quad L/F\text{ finite-unramified},\quad
 v\in S\Longrightarrow v\text{ splits in }L,\quad
 4\mid nd,\quad d>1.                                      \tag{CA1}
\]
The definite unitary group is compact at every real place
and quasi-split at every finite place.
Its existence condition is \(nd/2\) even, which is (CA1).
The fixed place \(v_0\) supplies the neatness condition of
Lemma~\ref{lem:ordinary-neat-level}; all auxiliary places
are outside \(S\).
One finite extension of the coefficient DVR \(\mathcal O\)
contains every required field embedding and every fixed
tame character value in (CA2). Further finite coefficient extensions
are taken faithfully flat.

At a split tame level place \(u\), put
\[
 \tau_u=\bigoplus_{i=0}^{n-1}\widetilde{\bar\chi_i}|_{I_u},
 \qquad
 \tau_0=\bigotimes_{u\in S,\ u\nmid\ell}
              \sigma_{\rm crys}(\tau_u)^\vee_{\mathcal O}.  \tag{CA2}
\]
Here the type lattice is the same finite free lattice as in
the ordinary function modules of
Section~\ref{sec:ordinary-module-construction}.
Away from the distinguished neat place \(v_0\), the
type theorem used in (CA2) is
Emerton--Gee, Theorem 9.2.4: its characteristic-zero occurrence
condition is the specified inertia type together with \(N=0\).
At \(v_0\), restriction of the chosen coefficient lattice
to the smaller neat subgroup does not impose the full
maximal-compact type. The independent \(N4\) argument
below proves Galois unramifiedness there.

Choose the local unit bases and the ordering of the residual
characters once. With
\[
 \Lambda=\mathcal O[[X_{vij}:
              v\mid\ell,\ 0\leq i<n,\ 1\leq j\leq d_v]],
\]
the character \(\theta^{-1}\) in the coefficient system is
characterized, in the reciprocity convention of the ordinary
modules, by
\[
 \psi_{vi}(g)\,\varepsilon(g)^i\,
       \theta(e_{vi}(y))=1
       \quad\text{if }\operatorname{Art}(g)=y\in\mathcal O_{L_w}^{\times}.
                                                               \tag{CA3}
\]
The normalized ordinary operators are
\[
 U_{vj}=\prod_{i=0}^{j}
              \psi_{vi}(g_v)\varepsilon(g_v)^i,
 \qquad \operatorname{Art}(g_v)=\varpi_v .
                                                               \tag{CA4}
\]
Thus both the unit values and the uniformizer values of
each \(\psi_{vi}\) are fixed. Equations (CA2)--(CA4), including
the dual in (CA2) and the inverse character in (CA3), are
the coefficient conditions required below.

\paragraph{Classical detectors with these coefficients.}
For a finite set \(Q\) of auxiliary split places, let
\(M_Q(K)\) be the residual-localized ordinary function module,
with coefficient lattice \(\tau_0\), character \(\theta^{-1}\),
and auxiliary compact subgroup
\(K\subset\prod_{v\in Q}\operatorname{GL}_n(L_w)\).
All other levels, including the neat level at \(v_0\), are fixed.
The constructions of
Lemma~\ref{lem:ordinary-projector-free} and
Proposition~\ref{prop:classical-ordinary-control} give finite
free \(\Lambda\)-modules and classical control for these
modules. The proof applies to smaller auxiliary compact
subgroups since their stabilizers still satisfy the same
neatness condition.

Here is the precise classical separation used with these
modules. Fix one sufficiently regular algebraic weight and
vary the finite-order characters of every local unit
coordinate. Require their \(n\) inertia characters at each
\(v\mid\ell\) to be pairwise distinct.
These arithmetic points are Zariski dense in \(\Lambda\):
a nonzero one-variable power series over \(\mathcal O\)
has only finitely many zeros among the characters under
consideration after Weierstrass preparation; induction on
the coordinates proves the assertion in several variables.
The forbidden equalities between two inertia characters
exclude only the corresponding proper character diagonals.
If \(z\ne0\) in \(M_Q(K)\), write \(z\) in a \(\Lambda\)-basis.
Some coefficient survives at one of these arithmetic points,
and classical control shows that the resulting classical
ordinary vector is nonzero. Decomposition of this classical
space into automorphic representations then detects \(z\).

For these detectors, the classical representation has a
continuous semisimple Galois representation with the
normalization
\[
 r_\ell(\Pi)=\operatorname{rec}
                   \bigl(\Pi|\det|^{(1-n)/2}\bigr).
                                                               \tag{CA5}
\]
This is used in the precise field and level scope (CA1).
One may apply Thorne,
\href{https://arxiv.org/pdf/1107.5989v1}{Theorem 6.5},
with its additional type set \(R\) empty; it allows the
remaining finite levels and does not require an irreducible
residual representation.
Equivalently, Labesse's Corollary 5.3, as stated in
\href{https://arxiv.org/pdf/0906.4189}{Guerberoff, Theorem 2.2},
gives the discrete isobaric base change, and Guerberoff's
Theorem 2.3 constructs its semisimple Galois representation.
The compact infinity factors satisfy Labesse's discrete-series
condition, the finite non-split factors are spherical, and
all prescribed ramified factors occur at split places.
If a discrete block is a residual Speh block, the algebraic
characters and powers of \(\varepsilon\) are combined as
in Section~\ref{subsec:sr-case-two-hecke}; this gives exactly
(CA5), rather than a different normalization on each block.

We use full local--global compatibility at auxiliary places.
For a regular algebraic cuspidal block this is
\href{https://annals.math.princeton.edu/wp-content/uploads/annals-v179-n2-p03-p.pdf}
{Barnet-Lamb--Gee--Geraghty--Taylor, Theorem 2.1.1}.
Taking the indicated direct sums and algebraic twists gives
the compatibility for each discrete isobaric base change
above. Thus its Weil--Deligne representation, including
monodromy, is a representation of the local Langlands
parameter used by the auxiliary calculation.
The weaker statement identifying only characteristic
polynomials would not suffice for this calculation.

The ordered ordinary flag at \(v\mid\ell\) also has the
characters in (CA3)--(CA4). Indeed, after killing the finite
inertia characters, local compatibility gives a crystalline
filtered module whose ordered Frobenius and descent root
lines are those selected by the ordinary Jacquet line.
For this assertion choose the fixed algebraic weight with
all successive labelled gaps sufficiently large.
Every cuspidal block then has Shin-regular weight, so
\href{https://smf.emath.fr/system/files/2017-08/ens_ann-sc_47_165-179.pdf}
{Barnet-Lamb--Gee--Geraghty--Taylor,
\emph{Local--global compatibility for \(\ell=p\), II},
Theorem 1.1} gives the Frobenius-semisimplified
local compatibility used here. The distinct finite
inertia characters give \(N=0\), so the finite extension
killing those characters makes this filtered module
crystalline.
For its first \(j\) root lines, the unit value of \(U_{v,j-1}\)
says that the Newton number is the sum of the \(j\) smallest
labelled Hodge--Tate weights. The induced filtration has
Hodge number at least that sum, whereas weak admissibility
gives the opposite inequality. Equality makes each of
these submodules weakly admissible. The corresponding
Galois subrepresentations form the full ordered flag.
The quotient characters are uniquely determined by their
unit values (CA3) and uniformizer values (CA4).
This is the rank-independent argument of
Section~\ref{subsec:sr-case-two-hecke}; it applies to the
ramified principal series just used for density.
An unramified principal-series ordinary-flag result alone
would not cover these detector points.

Let \(\mathbb T_Q(K)\) denote the selected ordinary Hecke
algebra acting faithfully on \(M_Q(K)\).
It is finite over \(\Lambda\).
The detectors just constructed are jointly faithful on
\(M_Q(K)\).
At each \(\ell\)-place the distinct finite inertial
characters make the Jacquet module of the corresponding
irreducible principal series the direct sum of its
Weyl-labelled character lines. Distinct character
lines are separated by a compact torus element;
the chosen coefficient character and ordering select
their indicated one-dimensional refinement.
Thus the selected ordinary operators act by scalars
on these lines, with no Jordan block. They act by
the same scalar on the accompanying automorphic
multiplicity space. The spherical operators act by
their spherical Hecke characters as well.
Consequently a nilpotent element of \(\mathbb T_Q(K)\)
vanishes on every detector and hence on \(M_Q(K)\);
faithfulness makes it zero. The algebra is reduced.
The integral determinant interpolation of
Lemma~\ref{lem:sr-compact-determinant-interpolation}
then gives a continuous determinant and local characters
over \(\mathbb T_Q(K)\), with (CA3)--(CA5).
For clarity, this interpolation uses integral characteristic
polynomial identities and compactness of the finite
\(\Lambda\)-algebra. It does not divide by \(n!\);
it remains valid when \(\ell\leq n\).
At the tame level, (CA2) imposes the stated inertial
relations. The polarization and ordinary determinant
relations hold on every classical detector and hence
over the Hecke algebra itself.

\paragraph{One ring acting on all auxiliary levels.}
Write \(\mathcal B_0\) for the base universal ordinary
determinant ring, \(D_0,\psi_0\) for its universal pair,
and
\[
 J_0=\ker(\mathcal B_0\longrightarrow\mathbb T_0),
 \qquad
 M_0=M_{\varnothing}.
\]
For the enlarged ramification set \(S\cup Q\), let
\(\mathcal R_Q\) be its universal ordinary determinant
ring, with its prescribed residual determinant and its
universal pair \(D_Q,\psi_Q\).
Thus \(\mathcal R_Q\), rather than
\(\mathbb T_Q(K)\), is the coefficient ring common to
all auxiliary compact subgroups \(K\).
The interpolation above and its closed relations give
\[
 f_K:\mathcal R_Q\longrightarrow\mathbb T_Q(K).
                                                               \tag{CA6}
\]
Before imposing the distinguished unramified relation,
apply the independent \(N4\) calculation leading to
(CA14) to the classical representations defining (CA6).
Classical separation then gives the determinant-kernel
relation itself, including its products with group-algebra
elements. This check uses the scalar residual
Frobenius polynomial and \(q_{v_0}\equiv-1\pmod\ell\);
it does not use a maximal-compact type at the smaller
neat level.
The universal property also gives the weight map
\(\lambda_Q:\Lambda\to\mathcal R_Q\). On the chosen
coordinates its formula is
\[
 \lambda_Q(X_{vij})=
       c_{vij}\,\psi_{Q,vi}(\gamma_{vj})-1,                 \tag{CA7}
\]
where \(c_{vij}\in\mathcal O^\times\) is the fixed
normalizing character value. In particular
\(f_K\lambda_Q\) is the scalar \(\Lambda\)-action, with
the same bases and constants for every \(Q,K\).

If \(K'\subset K\), the level inclusion
\(M_Q(K)\hookrightarrow M_Q(K')\) intertwines every
selected spherical and ordinary operator.
Restriction therefore gives
\(\mathbb T_Q(K')\to\mathbb T_Q(K)\).
Under this map the interpolated determinants agree:
their coefficients agree at the dense unramified
Frobenius classes, and continuity gives agreement
everywhere. Their local characters agree on units by
(CA3) and on a uniformizer by successive ratios of
(CA4); local reciprocity gives their equality on the
whole decomposition group. Uniqueness in the universal
ordinary determinant ring now gives compatibility of
the maps (CA6) with level inclusion.
The same argument applies to conjugated auxiliary levels.

Put
\[
 P_Q=\bigcup_K M_Q(K),\qquad
 G_Q=\prod_{v\in Q}\operatorname{GL}_n(L_w).
\]
Right translation makes \(P_Q\) a smooth
\(\Lambda[G_Q]\)-module. The compatible actions (CA6)
give a single homomorphism
\[
 \alpha_Q:\mathcal R_Q\longrightarrow
                  \operatorname{End}_\Lambda(P_Q),
 \quad
 \alpha_Q(\lambda_Q(a))z=az,\quad
 \alpha_Q(t)\rho(g)z=\rho(g)\alpha_Q(t)z .                 \tag{CA8}
\]
The last equality follows first for the generators of the
selected Hecke algebra and then for the universal ring
by the just-proved uniqueness. Thus (CA8) is a genuine
action on the smooth union, not a collection of unrelated
actions on finite modules.

Let \(h_Q:\Gamma_Q\to\Gamma_0\) be the natural map obtained
by imposing unramifiedness at the auxiliary places.
The base pair pulled back by \(h_Q\) gives
\[
 \beta_Q:\mathcal R_Q\longrightarrow\mathcal B_0,
 \quad
 D_Q\otimes_{\beta_Q}\mathcal B_0=D_0\circ h_Q,\quad
 \psi_Q\otimes_{\beta_Q}\mathcal B_0=\psi_0 .              \tag{CA9}
\]
At maximal auxiliary compact level there is an integral
identification
\[
 P_Q^{K_Q^{\max}}=M_0,\qquad
 \alpha_Q(t)|_{M_0}=\beta_Q(t)|_{M_0}.                    \tag{CA10}
\]
The first equality is the function-space fixed-vector
identity before localization. Localization preserves it:
a spherical auxiliary eigenpacket with the prescribed
residual determinant away \(S\cup Q\) has, by Chebotarev,
the same residual determinant everywhere that it is
unramified. Its missing auxiliary spherical residual
Hecke coefficients are therefore precisely the original
ones. Hence the two residual localizations agree.
For the second equality, compare the two Hecke actions on
\(M_0\). They agree on the Frobenius polynomials and,
by (CA3)--(CA4), on all local characters. The universal
property gives their equality. In particular the base
annihilator in (CA10) is always the single ideal \(J_0\).

For the Taylor--Wiles subgroups \(U_1\triangleleft U_0\)
at \(Q\), the earlier ordinary-module construction gives
\[
 P_Q^{U_1}\text{ finite over }\Lambda,\qquad
 P_Q^{U_1}\text{ projective over }\Lambda[U_0/U_1].
                                                               \tag{CA11}
\]
Together, (CA7)--(CA11) supply every ring, weight, finiteness,
projectivity and base-action field of the source
\srcid{LevMod}. Projectivity is the integral assertion
of Lemma~\ref{lem:ordinary-aux-projective}, not a deduction
from characteristic-zero dimensions.

\paragraph{The actual classical points of the smooth union.}
Fix an isomorphism
\(\iota:\overline{\mathbb Q}_\ell\simeq\mathbb C\).
For \(z\ne0\) in \(P_Q\), choose \(K\) with
\(z\in M_Q(K)\) and a classical detector at which \(z\)
survives. Choose an automorphic constituent and an ordered
ordinary refinement with nonzero projection of this
vector. Let \(\pi\) be its smooth auxiliary representation,
including its multiplicity space, and let
\[
 {\rm sp}:P_Q\longrightarrow Y
\]
be projection to this classical isotypic part after
specialization and extension through \(\iota\).
This map is additive and auxiliary-equivariant. Because
the action (CA8) was constructed from the same Hecke
operators, there are a scalar character
\(x:\mathcal R_Q\to\overline{\mathbb Q}_\ell\) and a
continuous semisimple Galois representation \(r_x\) such that
\[
 \begin{split}
 {\rm sp}(z)&\ne0,\\
 {\rm sp}(\alpha_Q(t)y)&=\iota(x(t))\,{\rm sp}(y),\\
 {\rm sp}(\rho(g)y)&=\pi(g){\rm sp}(y),\\
 D_Q\otimes_x\overline{\mathbb Q}_\ell&=\det r_x .
 \end{split}                                                \tag{CA12}
\]
For each \(v\in Q\), the restriction of \(Y\) to the local
factor is an isotypic sum of its irreducible local
constituent \(\pi_v\).
The Weil--Deligne compatibility already checked after
(CA5) is exactly the local compatibility field required
in \srcid{LevPtAlt}. Equations (CA12), for every nonzero
\(z\), prove \srcid{ClassicalAlt}.
They also explain why the same coefficient lattice and
the ramified ordinary flag are needed in this assertion.

We spell out the remaining arithmetic coefficient
identity in the old/new maps.
At an auxiliary place \(v\), let \(q_v\) be its residue
cardinality, \(\phi_v\) its chosen Frobenius, and
\(\Theta_v(t_i)\) the normalized torus operators on the
common unipotent-fixed module.
On its maximal-compact submodule,
\[
 e_k\bigl(q_v^0\Theta_v(t_0),\ldots,
          q_v^{n-1}\Theta_v(t_{n-1})\bigr)
       =D_Q^{[k]}(\phi_v)\quad(0\leq k\leq n).             \tag{CA13}
\]
Here \(D_Q^{[k]}\) denotes the signed characteristic
coefficient convention used by the determinant;
the \(q_v^i\) are the same normalizations as in the
old/new polynomial action.
To prove (CA13), apply every map \({\rm sp}\) in (CA12).
If the vector has zero image, both sides vanish on it.
Otherwise its maximal-compact invariance makes
\(\pi_v\) spherical. The Satake relation and local
compatibility identify the \(n\) normalized torus
characters with the roots of
\(\det(X-r_x(\phi_v))\). Vieta's formula gives (CA13)
over \(\mathbb C\).
Classical separation in (CA12) gives the integral
identity. Thus (CA13) is the arithmetic coefficient
comparison required by
Proposition~\ref{prop:concrete-two-maps}.

The geometric lemma, tame local description and finite
torus weight set now give the auxiliary relation ideal
with the uniform exponent of
Lemma~\ref{lem:torus-uniform-nil}.
For ranks two and three one may take \(n!\).
The intertwining maps are the actual two maps of
Proposition~\ref{prop:concrete-two-maps}, and their base
specializations are (CA10).
There is no change of \(J_0\) in passing from these
polynomial relations to the Taylor--Wiles modules.

At \(v_0\), the determinant in every classical detector
is unramified. Here one must use the scalar residual
Frobenius and
\(\bar\varepsilon(\phi_{v_0})=-1\) from the constructed
packet, together with the matrix \(N4\) argument.
Wild inertia has trivial residual action and maps
trivially to the pro-\(\ell\) congruence kernel.
For tame \(\ell\)-inertia the Frobenius conjugation
relation, scalar residual Frobenius, and
\(q_{v_0}\equiv-1\pmod\ell\) give the contraction
\(Y=T\,V^{-n}YQ\), with \(V\) invertible and \(T\)
topologically nilpotent, from the \(N4\) calculation.
Thus inertia is trivial. Detector separation gives
\[
 (D_0\bmod J_0)\text{ kills the closed inertia ideal at }v_0.
                                                               \tag{CA14}
\]
Neatness of the automorphic level alone would not prove
(CA14).

\begin{srproposition}[The arithmetic input for the common support calculation]
\label{prop:common-ordinary-actual-input}
For the constructed rank-two and rank-three packets,
the fixed ordinary Hecke ideal \(J_0\) admits the
Taylor--Wiles modules, comparison maps, coefficient
identities and uniform nilpotent relations on the
single base ring \(\mathcal B_0\) described in
(CA6)--(CA14). These data satisfy the exact
\srcid{HasAut} input of the common completed-support
calculation.
Consequently, using the arithmetic trace, obstruction
and prime-choice results of
Sections~\ref{sec:arithmetic-first-selmer},
\ref{sec:arithmetic-second-selmer} and
\ref{sec:arithmetic-tw-tower}, this same \(J_0\)
satisfies the characteristic-\(\ell\) conclusions
\srcid{P1core} and \srcid{P3core}.
\end{srproposition}

\begin{proof}
The first assertion is (CA6)--(CA13), (CA11) supplying
integral projectivity and (CA12) supplying the joint
classical detector family.
For any fixed finite list of nice curve points, normalize
their ordered eigenlines over finite coefficient
extensions and use the single Taylor--Wiles system of
Section~\ref{sec:arithmetic-tw-tower}.
Its coefficient-field extension can be chosen to have
one common finite cardinality for the entire list.
The trace detectors and obstruction estimates give
the two finite-level Selmer bounds for that same
system. The limit retains the original odd Selmer
generators and their stated spanning property; hence
Proposition~\ref{prop:tw-limit-coh} supplies its actual
limit cohomology and specialization maps.

Apply Proposition~\ref{prop:sr-one-patched-module} to
these modules. The completed support and its chain
are those of Proposition~\ref{prop:sr-completed-support-chain}.
At a type-I point use the ordinary Selmer presentation;
at a paired exceptional point use
Theorem~\ref{thm:selmer-paired-pres} and the completed
paired realization of Proposition~\ref{prop:paired-ch-realization}.
The arithmetic inclusion \((*)\) in that completed-support
chain is discharged in
Section~\ref{subsec:sr-same-completion-galout}:
the one added formal weight variable maps to the selected
curve parameter, so the regular weight completion has
dimension \(1+nd\). The \(n|Q|\) auxiliary variables
give support dimension \(1+nd+n|Q|=1+D_0\), exactly the
bound \(D_0=nd+n|Q|\) in the first-Selmer calculation.
The coefficient DVR section and the first-/second-order
tests are those in this same completion.
Lemma~\ref{lem:sr-ordinary-presentation-prime} gives the
single-block inclusion; the odd-product detector,
Corollary~\ref{cor:paired-nonss}, and
Lemma~\ref{lem:sr-paired-presentation-prime} give the
paired inclusion from the actual Selmer presentation.
Thus the application does not assume \((*)\).
The cross-cocycle finiteness and the closed
Cayley--Hamilton quotient required there are supplied
by Proposition~\ref{prop:sr-finite-CH}.
The closure is the joint profinite-group and adic
closure used in that proposition.

All modules in this step are the modules of (CA8);
the distinguished maximal-compact fiber is \(M_0\)
by (CA10). Thus annihilator transport and coefficient
descent return the original \(J_0\).
Faithful flatness of the finite coefficient extension
allows descent to the original \(\mathcal O\).
The completed support chain gives the two assertions:
at a type-I curve, every specialization of its kernel
under consideration lies in the support of \(J_0\);
for a finite chain of exceptional curves, the same
holds with the specified nonsemisimple closed
Cayley--Hamilton quotients. These are precisely
\srcid{P1core} and \srcid{P3core}.
No residual Bernoulli-unit or residual unramified
cohomology-vanishing assertion enters this construction.
\end{proof}

\paragraph{Exact source trace.}
The finite-level coefficient construction above
reconstructs
\srcid{TCoreSep.coreSep\_coord} and
\srcid{TCoreSep.coreSep\_kill},
\srcid{TCoreDet.coreDet\_alt},
\srcid{LevelUR2Compat.level\_change\_compat\_rank\_two\_of\_dense\_weight},
\srcid{LevelUR2ActEx.levelUR2ActEx\_alt},
\srcid{LevelUR2BaseEq.levelUR2BaseEq\_alt},
\srcid{LevelUR2SmulKey.levelUR2SmulKey\_alt}, and
\srcid{NonemptyUnitaryHeckeParts.levelUR2ModPack\_alt}
under
\srcid{ModularityArgument.ThresholdArguments.DenseThresholdCases}.
Despite some source names, their displayed statements
are rank-\(n\) statements.
The coefficient identity (CA13) is
\srcid{SymOfVieta.symOf\_alt};
the two-map and uniform-relation assembly is
\srcid{THasAutOf.has\_automorphy\_of\_statement}.
The final common support descent is the actual chain
\srcid{SeparableMatrixStability.LiftPt.core}
\(\longrightarrow\)
\srcid{CoreBC}
\(\longrightarrow\)
\srcid{RawMain.raw\_main}
\(\longrightarrow\)
\srcid{ChebotarevFamilies.PFromAutOf.chain\_of}
and \srcid{pFromAut\_of}.
These source names specify provenance; the construction
and comparison needed at each arrow have been given above.

\subsection{The same completed ring in the local support inclusion}
\label{subsec:sr-same-completion-galout}

We spell out the local inclusion used in
Proposition~\ref{prop:common-ordinary-actual-input}.
This also constructs the coefficient section required in
Theorem~\ref{thm:selmer-paired-pres}.

At a selected normalized curve, extend the coefficient DVR once
so that its residue field is the curve's finite field \(k\).
Write \(A_0=k[[T]]\), \(F=k((T))\).
Let \(R\) be the single enlarged coefficient ring of the common
patched module, let \(I\) be its already imposed relation ideal,
and let
\[
 B_0=R/I,\qquad f:B_0\twoheadrightarrow A_0,\qquad
 B=\widehat{(B_0)_{\ker f}} .
\]
The point supplies \(s\in B_0\) with \(f(s)=T\).
All determinants and characters below are the pair over this
\(B_0\), extended to this \(B\).
Neither \(B\) nor \(I\) is replaced by the image algebra of an
auxiliary automorphic module.

\paragraph{An explicit compatible coefficient DVR.}
Choose a uniformizer \(\pi\) of \(\mathcal O\).
The continuous evaluation \(Z\mapsto s\) defines
\(\mathcal O[[Z]]\to B_0\).
Its composite with \(f\) is coefficient reduction followed by
\(Z\mapsto T\); its kernel is exactly \((\pi)\).
The residue field of \(B\) is \(F\):
localization inverts every nonzero element of \(A_0\), and the
surjectivity of \(f\) supplies every such fraction.
Consequently the inverse image of \(\mathfrak m_B\) in
\(\mathcal O[[Z]]\) is \((\pi)\).
Evaluation therefore extends to a local map
\[
 \mathcal O_0=
    \widehat{\mathcal O[[Z]]_{(\pi)}}\longrightarrow B .
 \tag{LCJ1}
\]
The localization is a DVR, since \(\mathcal O[[Z]]\) is regular
and \((\pi)\) is a height-one prime.
Its completion \(\mathcal O_0\) is a complete DVR with residue
field \(F\), and (LCJ1) induces the identity on \(F\).
Extension of the map uses completeness of \(B\);
a \(\pi\)-adic Cauchy sequence maps to a
\(\mathfrak m_B\)-adic Cauchy sequence.

This section has all three required compatibilities.
Coefficientwise lifts in \(\mathcal O[[Z]]\) of an arbitrary
element of \(k[[T]]\) give a common lift in \(\mathcal O_0\) and
\(B_0\); constants give the original \(\mathcal O\)-lifts;
and \(Z\) maps to the specified element \(s\).
Thus the section assumptions in the finite-test boundedness
argument of Theorem~\ref{thm:selmer-paired-pres} hold here
without an additional coefficient-field assumption.
This is the explicit construction in
\srcid{SelmerScaling.OoExKer.ooEx}.

\paragraph{The dimension is the same integer as in the Selmer bound.}
Adjoin the one formal variable \(S\) to the weight ring and to
all common modules before patching; the curve map sends \(S\)
to \(T\).
This formal extension preserves freeness, the two control
maps, their uniform annihilator relations, and the extended
base ideal \(J_0\).
The augmented weight ring is
\[
 \Lambda^\dagger=\mathcal O[[S,X_1,\ldots,X_{nd}]] .
\]
It has dimension \(nd+2\).
Its map to \(k[[T]]\) is surjective because \(S\mapsto T\),
so the quotient by its kernel has dimension one.
Regularity and the dimension formula show that completion at
this kernel is regular of dimension \(nd+1\).
If the Taylor--Wiles system has \(q_{\rm TW}\) places, its
\(nq_{\rm TW}\) auxiliary variables consequently give
\[
 \dim(\text{each completed support component})
      =nd+1+nq_{\rm TW}=D_0+1,\qquad
 D_0=nd+nq_{\rm TW}.
 \tag{LCJ2}
\]
Here the first equality is precisely
Lemma~\ref{lem:sr-patched-support-dimension}.
The integer \(D_0\) on the right is precisely the first-Selmer
bound for this same Taylor--Wiles system, as supplied by
Proposition~\ref{prop:tw-limit-coh}.
This is the dimension matching of
\srcid{AnnihilatorBaseChange.LaDimArith.laDim}.
In particular the dimension threshold is not obtained by
completing the unaugmented weight ring and discarding its
one-dimensional curve quotient.

\paragraph{The actual local inclusion.}
The finite closed Cayley--Hamilton construction and
Proposition~\ref{prop:paired-ch-realization}, applied over
\(B_0\) and then extended to \(B\), give the algebra
\(\mathcal A=B\oplus N\) and its actual representation.
The point's block span is the residue span used in that
proposition.  Its two cases are the one-block case and the
paired-block case.
The finite group/adic closure was imposed in \(B_0\) before
scalar extension; hence the kernel comparison is the one
required in Corollary~\ref{cor:paired-nonss}.

Take primes \(Z,Q\subset B\) occurring in the completed
support chain, with \(Z\) stable and
\(\dim B/Z\ge D_0+1\), and suppose that \(Q\) is stable or
has the specified nonsemisimple closed Cayley--Hamilton
quotient.
The actual limit cohomology gives the three Laurent inputs
of Theorem~\ref{thm:selmer-paired-pres}, with bound \(D_0\).
The section (LCJ1) and the finite-lattice continuity already
proved for the closed algebra verify its remaining hypotheses.

In the one-block case \(N=0\), that theorem gives
\[
 \mathfrak m_B\subset
   (X_1,\ldots,X_{D_0})+\mathfrak m_{\mathcal O_0}B+
                                      \mathfrak m_B^2 .
\]
Lemma~\ref{lem:sr-ordinary-presentation-prime} gives \(Z=0\),
hence \(Z\subseteq Q\).

In the paired-block case put \(I_{\rm odd}=(NN)\subset B\).
Lemma~\ref{lem:sr-odd-product-detector} makes stability at
\(Z\) imply \(I_{\rm odd}\not\subseteq Z\).
At \(Q\), stability gives \(I_{\rm odd}\not\subseteq Q\);
the nonsemisimple alternative instead gives
\(\operatorname{Ann}_B N\subseteq Q\), by
Corollary~\ref{cor:paired-nonss}.
Theorem~\ref{thm:selmer-paired-pres} gives the independent
coefficient rows and the relations required by
Lemma~\ref{lem:sr-paired-presentation-prime}.
That lemma gives \(Z\subseteq Q\) in both alternatives.

This proves exactly the local inclusion \((*)\) of
Proposition~\ref{prop:sr-completed-support-chain}.
Its stable and nonsemisimple predicates are those of the
same extended determinant, and (LCJ2) gives its actual
dimension threshold.
All its hypotheses now refer to the single ring \(R\), the
single relation ideal \(I\), and the same common module.
Thus the chain's annihilator transport returns the original
base Hecke ideal \(J_0\).
Together with (CA6)--(CA14), this supplies the local
component-inclusion step used for \srcid{P1core} and
\srcid{P3core}; it is not a further automorphy hypothesis.
